% Modernised reading — fonds Alexandre Grothendieck, shelfmark 56, pages 1-38.
%
% This is an INTERPRETATION, not a transcription. Everything the critical
% apparatus carried moves into footnotes. In any dispute of substance the
% transcription governs: whoever cites Grothendieck must cite batch-01.fr.tex
% and batch-02.fr.tex.
%
% Pass: Opus 5.5 (claude-opus-5-5), 2026-10-03 — first pass, unchecked against the pages by a human.
% The folder was read whole: two batches, pages 1-20 and 21-38, both
% transcribed. Blank: 9, 18, 33; covers: 1, 8; colophon only: 38.
%
% What the folder is. Three things in one sleeve. (1) Manuscript notes on the
% vanishing cycles of a degeneration of a smooth quadric to a quadric with
% one ordinary double point (pages 2, 4, 6), and the computations they rest
% on: cohomology of a projective cone and of its punctured affine cone
% (10, 12, 14, 16), of a smooth quadric and its affine part (19, 21), the
% cohomology ring of quadrics and its functoriality (23, 25, 27, 29, 31), a
% fair note on their Chow groups (34). (2) Fourteen leaves of his own
% typescript of SGA 3 Exp. XII (Tores maximaux, groupe de Weyl...), corrected
% in ink, on whose backs the notes are written (odd pages 3-17, 20, even
% pages 22-32). (3) An offprint of A. Delzant's CRAS Note of 1962 on
% Stiefel-Whitney classes of quadratic forms (35-37), with nothing of his.
%
% Notation fixed for the whole document. Coefficients Λ = Z/m, m invertible
% on the base, Λ(i) the Tate twist (the pages write μ(i), μ_m(i)); for a
% Λ-module A, A(i). Pages 2-6: S strictly local trait, s, η, \bar η; X_s and
% X_{\bar η} for his X_0 and \bar X_1; fibres of dimension n-1 (his n);
% I the inertia (his π). Vanishing cycles in the modern index:
% H^i(X_s, RΦΛ) is his Φ^{i+1}. Pages 10-16: the base of the cone is
% called Y (the page: X) and its dimension d (the page: n), to keep X and
% n for pages 2-6. Pages 19-34: Q^n a smooth quadric of dimension n, U^n its
% affine part. Classes: ξ hyperplane; ξ_i the class of a linear subspace of
% codimension i for i > n/2 (his « ½ξ^i »); on Q^{2ν}, α_ν, β_ν the classes
% of the two families of ν-planes, λ_ν = α_ν - β_ν, φ_ν = α_ν; ℓ^i, ℓ_i
% generators of H^i(U^i) and H^i_c(U^i).
%
% Substantive departures from the pages, each footnoted where it occurs:
% — pages 4, 12, 23: [ξ^ν, λ_ν] is a basis only once 2 is invertible; the
%   integral basis is (ξ^ν, φ_ν), as page 23's own margin says. The cokernel
%   of page 4 (n odd) is generated by φ_ν, not λ_ν, when 2 is not invertible.
% — page 4, n odd: the action of inertia on the vanishing cycle is the
%   quadratic character by which it exchanges the two families; nontrivial
%   when X is regular (residue characteristic ≠ 2). The page's « non » is
%   uncertain.
% — page 6: the last two terms of the sequence are for U, not X; the
%   hypothesis that Z is smooth over S and misses the singular point is
%   supplied.
% — page 12: for i ≥ d+1 the right exponent is i-1 (kernel of ξ on
%   H^{i-1}), not i-2; hard Lefschetz is what is used.
% — page 14: torsion sits in degree 2ν, not 2ν-2 (first case); the second
%   case has H^{2ν}(E) ≅ Λ(-ν), not H^{2ν-1}(E) ≅ Λ(-ν+1).
% — page 19: ν'_n = Λ(-(n-1)/2) for n odd, not -(n+1)/2 (the added line in
%   another ink is right); ν_n = Λ(-(n+1)/2) for n odd; ⟨ℓ, φ(ℓ)⟩ = (-1)^{n/2}·2.
% — page 21: c), c') stated in their true form; d) « i = n » read as 2i = n.
% — page 23: Gysin raises the index by 1, not 2 (page 25 has it right).
% — page 25: λ_ν^2 = (-1)^ν·2ξ_{2ν}, not 2ξ_{2ν} for all ν.
% — page 29: the first term of the upper row read H^{n+1}(U^{n+1})(1) (the
%   page: H^{n-1}(U^{n-1})(1)); the « (+1) » of the lower row dropped; the
%   « bij » of the curved arrow is multiplication by 2.
%
% Modern names given and footnoted as ours: cycles évanescents (SGA 7),
% formule de Picard-Lefschetz, point double ordinaire, entrelacs, théorème de
% Lefschetz difficile, cohomologie primitive, lemme de contraction, anneau de
% Chow des quadriques, théorème de Merkurjev (for the offprint's questions).
%
% What the folder announces and never establishes: the boxed
% « H*_a(C) ≅ H*_c(C) ?? » of page 16 (true, by the contraction lemma); the
% e) of page 21 (illegible); the sense of « pur » on page 4.
\documentclass[11pt,a4paper]{article}
\input{../preamble/grothendieck}

\folder{56}
\pages{1}{38}
\dating{1962}
\watermark{Édition de démonstration\\interprétation personnelle de l'œuvre}
\foldertitle{Lefschetz-Hodge - Vanishing cycles : notes manuscrites (s.d.), tiré à part (1962) — lecture modernisée du dossier entier}

\begin{document}

\section*{Résumé}

\begin{resume}
Ces feuillets étudient ce qui arrive à la forme d'une surface, ou d'un objet
de dimension plus grande, au moment où une famille d'objets lisses se
pince en un point. L'exemple de base tient dans une équation :
$x^{2} + y^{2} + z^{2} = t$. Pour $t$ positif, les points réels forment une
sphère ; quand $t$ tend vers zéro, la sphère rétrécit et disparaît dans un
point singulier, le sommet du cône $x^{2} + y^{2} + z^{2} = 0$. La sphère
qui s'évanouit s'appelle un \emph{cycle évanescent}. Elle porte une classe de
cohomologie de la fibre lisse qui n'a plus de place dans la fibre singulière,
et c'est elle qui mesure la différence entre les deux.

Le problème est de rendre ce tableau exact en cohomologie étale, où il n'y a
plus de sphère à dessiner, et de le faire sur l'exemple le plus simple où il
dit quelque chose : une quadrique lisse qui dégénère en une quadrique à un
seul point double, le cône sur une quadrique de dimension deux de moins.
Les notes comparent la cohomologie des deux fibres, isolent la différence
dans des groupes de cycles évanescents, et trouvent qu'elle tient en un seul
degré, celui de la dimension de la fibre, et en un seul générateur. Elles
regardent aussi comment ce générateur se comporte quand on tourne autour du
point où la famille dégénère : selon la parité de la dimension, il revient
à lui-même ou à son opposé. C'est, sur un exemple, ce qu'on appelle
aujourd'hui la formule de Picard-Lefschetz.

Pour faire ce calcul il faut connaître la cohomologie de tous les
morceaux, et c'est ce que fait l'essentiel du dossier : la cohomologie d'un
cône projectif et du cône épointé, celle d'une quadrique lisse et de sa
partie affine, l'anneau de cohomologie des quadriques avec une base
explicite. Un nombre y joue un rôle à part, le nombre $2$ : une quadrique
est de degré deux, et la classe d'un sous-espace linéaire y vaut la moitié
d'une puissance de la section hyperplane. Les bases écrites « à $2$
près » dans les premières pages sont rectifiées dans les suivantes, et
c'est encore un facteur $2$ qui sépare le cycle évanescent de son image dans
la fibre singulière. La lecture qui suit donne les énoncés sous leur forme
exacte, y compris là où le brouillon se trompe d'un degré ou d'un signe.

Les notes sont écrites au dos d'un tapuscrit corrigé de sa main, celui de
l'exposé XII du séminaire SGA~3 sur les tores maximaux des schémas en
groupes. On y voit l'auteur réviser son propre texte, et noter en marge que
le contenu de quatre exposés « devrait être complètement refondu ». Le
dossier contient enfin le tiré à part d'une Note aux Comptes rendus de
1962 sur les classes de Stiefel-Whitney des formes quadratiques, qui donne
à l'inventaire sa date ; les notes, elles, ne sont pas datées.

\keywords{vanishing cycles, Picard-Lefschetz formula, ordinary double point, quadric, projective cone, Gysin sequence, local cohomology, link of a singularity, hard Lefschetz theorem, primitive cohomology, contraction lemma, Chow ring of quadrics, maximal torus, Weyl group, reductive rank, Stiefel-Whitney classes of quadratic forms}
\end{resume}

\section*{Le fil du dossier, et les conventions}

Les notes manuscrites, au dos des feuillets du tapuscrit, se suivent sous
deux chemises qui portent leur titre de sa main. Elles font un seul
calcul, pris par les deux bouts :

\begin{enumerate}
\item[1.] \emph{La dégénérescence} (pages 1 à 6) : la quadrique lisse, sa
dégénérée, la suite exacte de spécialisation, le résultat dans les deux
parités, puis le même résultat par un « calcul affine » plus direct.
\item[2.] \emph{Le cône projectif} (pages 8 à 16) : la cohomologie du cône,
du cône épointé, la cohomologie locale au sommet, le cas d'un cône sur une
quadrique, les supports compacts. La fibre singulière du point 1 est un tel
cône.
\item[3.] \emph{La quadrique lisse et sa partie affine} (pages 19 et 21) :
la partie affine est la fibre lisse du « calcul affine » du point 1.
\item[4.] \emph{L'anneau de cohomologie des quadriques} (pages 23 à 34) :
bases, relations, restrictions et Gysin entre quadriques emboîtées, le carré
de la page 29 qui relie le cycle évanescent à la classe $\lambda_{\nu}$, le
cup-produit par $\xi$, et une note au propre sur les groupes de Chow.
\end{enumerate}

Les quatorze feuillets du tapuscrit sont lus à part, après les notes, dans
l'ordre de leurs paragraphes ; le tiré à part vient en dernier.

\emph{Conventions.} Les coefficients sont $\Lambda = \mathbf{Z}/m$, $m$
inversible sur la base, et $\Lambda(i)$ est le tordu de Tate ; pour un
$\Lambda$-module $A$, on écrit $A(i)$.\footnote{Les pages écrivent
$\mu(i)$, avec un $\mu$ à jambage doublé, pour ce que nous notons
$\Lambda(i)$, et $\mu_{m}(i)$ à la page 19. Tout ce qui suit vaut aussi, par
passage à la limite, pour $\mathbf{Z}_{\ell}$.} Deux lettres servent deux
fois dans les pages, et on les sépare ici. Aux pages 2 à 6, $X$ est
l'espace total de la famille et ses fibres sont de dimension $n - 1$,
comme sur la page. Aux pages 10 à 16, la base du cône, que la page appelle
aussi $X$, est notée $Y$, et sa dimension, que la page appelle $n$, est
notée $d$. À partir de la page 19, $Q^{n}$ est une quadrique lisse de
dimension $n$, $Q^{n-1}$ une section hyperplane lisse et
$U^{n} = Q^{n} - Q^{n-1}$ la quadrique affine.

Sur une quadrique, $\xi$ est la classe de la section hyperplane. Pour
$i > n/2$, la classe d'un sous-espace linéaire de codimension $i$ contenu
dans $Q^{n}$ est notée $\xi_{i}$ ; on a $2\xi_{i} = \xi^{i}$, et la page
l'écrit $\tfrac12 \xi^{i}$, notation qui n'est littérale que si $2$ est
inversible. Sur $Q^{2\nu}$, les $\nu$-plans contenus dans la quadrique
forment deux familles ; on note $\alpha_{\nu}$, $\beta_{\nu}$ leurs classes,
$\lambda_{\nu} = \alpha_{\nu} - \beta_{\nu}$ et
$\varphi_{\nu} = \alpha_{\nu}$, de sorte que
$\xi^{\nu} = \alpha_{\nu} + \beta_{\nu}$ et
$\varphi_{\nu} = \tfrac12(\xi^{\nu} + \lambda_{\nu})$ : ce sont les
définitions de la page 29.

\emph{Ce que le dossier annonce et n'établit pas} : l'égalité encadrée
« $H^{*}_{a}(C) \simeq H^{*}_{c}(C)$ ?? » de la page 16, qui est vraie ;
le point e) du programme de la page 21, illisible ; le sens du mot
« pur » à la page 4.

\pagerange{1}{2}
\section*{La dégénérescence d'une quadrique (pages 1 à 6)}

\subsection*{La situation (page 2)}

Soit $S$ le spectre d'un anneau de valuation discrète strictement
hensélien, de point fermé $s$, de point générique $\eta$, et $\bar\eta$ un
point géométrique au-dessus de $\eta$ ; soit $I$ le groupe d'inertie,
qui est ici le groupe de Galois de $\bar\eta$ sur $\eta$.\footnote{La page
écrit « strictement local » et note $\pi$ ce groupe ; elle note $X_{0}$
et $\overline{X}_{1}$ les fibres que nous notons $X_{s}$ et
$X_{\bar\eta}$. La chemise de la page 1 porte « Situation des
« vanishing cycles » : [un mot illisible] Cas des quadriques », et
le coin de la page 2 répète « Vanishing cycles : cas des quadriques ».}
Soit $X \subset \mathbf{P}^{n}_{S}$ défini par une équation quadratique,
$f : X \to S$ projectif et plat, de fibres de dimension $n - 1$, tel que
la fibre générique géométrique $X_{\bar\eta}$ soit une quadrique lisse et
que la fibre spéciale $X_{s}$ soit une quadrique avec un seul point
singulier $x$. Une telle quadrique est le cône projectif, de sommet $x$,
sur une quadrique lisse de dimension $n - 2$, et $x$ est un point double
ordinaire.\footnote{« Point double ordinaire » est notre mot ; la page
dit « quadrique à point singulier ». Une quadrique de $\mathbf{P}^{n}$
sur un corps algébriquement clos n'a qu'un point singulier exactement
quand sa forme est de rang $n$ : c'est alors le cône sur une quadrique
lisse de $\mathbf{P}^{n-1}$.}

La cohomologie de $X_{s}$ s'identifie à celle de $X$ (changement de base
propre, $S$ étant strictement local), et la restriction à $X_{\bar\eta}$
donne la flèche de cospécialisation
$H^{i}(X_{s}) \to H^{i}(X_{\bar\eta})$, compatible à l'action de $I$, qui
agit trivialement sur la source. La question de la page est de comprendre
cette flèche. On l'insère dans la suite exacte
\[
  \cdots \to H^{i}(X_{s}) \to H^{i}(X_{\bar\eta}) \to H^{i}(X_{s}, R\Phi\Lambda)
  \to H^{i+1}(X_{s}) \to H^{i+1}(X_{\bar\eta}) \to \cdots
\]
qui naît du triangle $\Lambda \to R\Psi\Lambda \to R\Phi\Lambda \to$ des
cycles proches et évanescents.\footnote{Le formalisme $R\Psi$, $R\Phi$ est
celui de SGA~7 (exposé XIII, 1967-1969). Les notes ne l'ont pas : elles
écrivent la suite avec des groupes $\Phi^{i}$ qu'elles définissent par elle,
comme « les noyaux et conoyaux » de la spécialisation, et leur
indexation est celle d'une cohomologie relative : leur $\Phi^{i+1}$ est
notre $H^{i}(X_{s}, R\Phi\Lambda)$. Nous suivons l'indexation moderne, où le
groupe intéressant est en degré $n - 1$, la dimension des fibres, et non en
degré $n$ comme sur la page.} La classe $\xi \in H^{2}(X, \Lambda(1))$ de
$\mathcal{O}_{X}(1)$ se restreint en les classes hyperplanes des deux
fibres : c'est ce qui fera correspondre les puissances de $\xi$.\footnote{Le
N.B. qui termine la page est d'une écriture très rapide ; seuls $X$, $X_{0}$
et la classe $\xi \in H^{2}(X, \mu(1))$, « provenant de la section »,
se lisent sûrement. La lecture du N.B. comme restriction de la classe
hyperplane est la nôtre.}

Les deux fibres n'ont pas de cohomologie en degré impair (on le verra pour
le cône aux pages 10 et 12, et pour la quadrique lisse à la page 21). La
suite se coupe donc en morceaux
\[
  0 \to H^{2k-1}(X_{s}, R\Phi\Lambda) \to H^{2k}(X_{s}) \to H^{2k}(X_{\bar\eta})
  \to H^{2k}(X_{s}, R\Phi\Lambda) \to 0 ,
\]
et les groupes de cycles évanescents sont les noyaux et les conoyaux de la
cospécialisation en degré pair.\footnote{La page écrit ce morceau sous
« 1°) $n$ pair », avec « $n$ impair » biffé à côté ; l'argument
vaut dans les deux parités, puisqu'il n'utilise que l'annulation en degré
impair.}

\pagerange{4}{4}
\subsection*{Le calcul dans les deux parités (page 4)}

La fibre lisse $X_{\bar\eta}$ est une quadrique de dimension $n - 1$, la
fibre spéciale le cône sur une quadrique de dimension $n - 2$ ; les bases
de leur cohomologie viennent des pages 12, 23 et 27.

\emph{$n = 2\nu$ pair.} Les fibres sont de dimension impaire $2\nu - 1$.
\[
  \begin{array}{ll}
    H^{*}(X_{s}) : & 1, \xi, \ldots, \xi^{\nu-1}, \ [\xi^{\nu}, \varphi'_{\nu}], \
      \xi_{\nu+1}, \ldots, \xi_{2\nu-1} \\[4pt]
    H^{*}(X_{\bar\eta}) : & 1, \xi, \ldots, \xi^{\nu-1}, \ \xi_{\nu}, \
      \xi_{\nu+1}, \ldots, \xi_{2\nu-1}
  \end{array}
\]
Dans $X_{s}$, $\varphi'_{\nu}$ est la classe d'un $(\nu-1)$-plan de
l'une des deux familles de la quadrique de base, vue dans le cône (c'est
$j_{*}\varphi_{\nu-1}$ avec les notations de la page 10), et
$\lambda'_{\nu} = 2\varphi'_{\nu} - \xi^{\nu} = j_{*}\lambda_{\nu-1}$ la
différence de deux tels plans de familles opposées. La cospécialisation envoie $\xi$ sur $\xi$,
$\varphi'_{\nu}$ sur la classe $\xi_{\nu}$ d'un $(\nu-1)$-plan de
$X_{\bar\eta}$, donc $\lambda'_{\nu}$ sur $0$.\footnote{La page écrit la
base $[\xi_{0}^{\nu}, \lambda'_{\nu}]$ ; elle n'est une base que si $2$ est
inversible dans $\Lambda$, et avec elle la surjectivité qui suit serait
fausse pour $\ell = 2$, puisque l'image de $\xi^{\nu}$ est $2\xi_{\nu}$. La
base entière $(\xi^{\nu}, \varphi'_{\nu})$ est celle que la page 23 donne en
marge pour la quadrique lisse ; nous la prenons ici.} Elle est donc bijective
en tout degré $i \neq n$, surjective en degré $n = 2\nu$, de noyau libre de
rang un engendré par $\lambda'_{\nu}$. D'où :

\begin{enumerate}
\item[a)] $H^{i}(X_{s}, R\Phi\Lambda) = 0$ pour $i \neq n - 1$, et
$H^{n-1}(X_{s}, R\Phi\Lambda) \simeq \Lambda(-\nu)$, l'isomorphisme
n'étant déterminé qu'au signe près (le choix de l'une des deux
familles) ;\footnote{« isom.\ déterminé au signe près » est de la page ;
l'exposant de $\Phi$ et le $-\nu$ y sont repassés et empâtés.}
\item[b)] $R^{i}f_{*}\Lambda$ est constant pour $i \neq n$, et
$R^{n}f_{*}\Lambda$ est extension d'un faisceau constant par un faisceau
concentré au point fermé $s$, isomorphe à $\Lambda(-\nu)$ ;\footnote{La page
dit « $R^{n}f_{*}$ = un [mot illisible]-faisceau concentré en $s$, isom.\
à $\mu(\ldots)$ », le twist laissé en blanc. Un faisceau sur le trait
strictement local $S$ est donné par sa fibre en $s$, sa fibre en $\bar\eta$
avec l'action de $I$, et la flèche de cospécialisation ; ici la flèche est
surjective et $I$ agit trivialement, de sorte que $R^{n}f_{*}\Lambda$ contient
le faisceau gratte-ciel de valeur $\Lambda(-\nu)$ en $s$, de quotient
constant. C'est notre lecture du mot manquant.}
\item[c)] $I$ agit trivialement sur $H^{*}(X_{\bar\eta})$ et sur les cycles
évanescents.\footnote{En effet le groupe des cycles évanescents s'injecte
dans $H^{n}(X_{s})$, où $I$ agit trivialement. Le point c) est écrit dans la
marge, en regard de b).}
\end{enumerate}

\emph{$n = 2\nu + 1$ impair.} Les fibres sont de dimension paire $2\nu$.
\[
  \begin{array}{ll}
    H^{*}(X_{s}) : & 1, \xi, \ldots, \xi^{\nu}, \ \xi_{\nu+1}, \ldots, \xi_{2\nu} \\[4pt]
    H^{*}(X_{\bar\eta}) : & 1, \xi, \ldots, \xi^{\nu-1}, \ [\xi^{\nu}, \varphi_{\nu}], \
      \xi_{\nu+1}, \ldots, \xi_{2\nu}
  \end{array}
\]
La cospécialisation est bijective pour $i \neq n - 1 = 2\nu$ ; en degré
$2\nu$ elle envoie $\xi^{\nu}$ sur $\xi^{\nu} = \alpha_{\nu} + \beta_{\nu}$ :
elle est injective, de conoyau libre de rang un engendré par l'image de
$\varphi_{\nu} = \alpha_{\nu}$.\footnote{La page dit le conoyau engendré par
$\lambda_{\nu}$, le symbole surchargé au point d'être à peine lisible. Or
$\lambda_{\nu} = 2\alpha_{\nu} - \xi^{\nu}$ n'a pour image que le double du
générateur ; il n'engendre le conoyau que si $2$ est inversible.} D'où :

\begin{enumerate}
\item[a)] $H^{i}(X_{s}, R\Phi\Lambda) = 0$ pour $i \neq n - 1$, et
$H^{n-1}(X_{s}, R\Phi\Lambda) \simeq \Lambda(-\nu)$ comme groupe ;
\item[b)] $R^{i}f_{*}\Lambda$ est constant pour $i \neq n - 1$, et la
cospécialisation de $R^{n-1}f_{*}\Lambda$ est injective : ce faisceau n'a
pas de section non nulle concentrée en $s$ ;\footnote{La page dit
« $R^{n-1}f_{*}$ est « pur »  » et s'arrête. Le sens donné ici au
mot, l'absence de sous-faisceau concentré au point fermé, est notre
interprétation, la seule que les calculs de la page appuient.}
\item[c)] $I$ agit trivialement sur $H^{i}(X_{\bar\eta})$ pour
$i \neq n - 1$ ; sur $H^{2\nu}(X_{\bar\eta})$ il permute les deux familles
de $\nu$-plans, donc agit sur les cycles évanescents par un caractère
$\varepsilon$ d'ordre au plus $2$, et l'on a
$H^{n-1}(X_{s}, R\Phi\Lambda) \simeq \Lambda(-\nu) \otimes \varepsilon$. Si
la caractéristique résiduelle n'est pas $2$ et si $X$ est régulier,
$\varepsilon$ n'est pas trivial.\footnote{La page écrit « mais
non sur $\Phi^{n}$ », le « non » lu avec doute, dans la marge
au pied de la page. La précision est la nôtre : en caractéristique
résiduelle différente de $2$, on diagonalise la forme en
$x_{0}^{2} + \cdots + x_{n-1}^{2} + c\,x_{n}^{2}$ avec $c$ dans l'idéal
maximal ; les deux familles de $X_{\bar\eta}$ sont définies sur
$\eta(\sqrt{\pm c})$, ramifiée exactement si la valuation de $c$ est
impaire, et $X$ est régulier au sommet exactement si cette valuation vaut
$1$. Avec le cas pair, c'est la formule de Picard-Lefschetz pour une
singularité quadratique, établie en général par Deligne dans SGA~7,
exposé XV : action triviale sur les cycles évanescents en dimension
relative impaire, par un caractère quadratique en dimension relative
paire.}
\end{enumerate}

\pagerange{6}{6}
\subsection*{Le calcul affine (page 6)}

Le même résultat s'obtient directement, sans connaître la cohomologie des
fibres projectives. Soit $Z = X \cap H$ une section hyperplane qui ne passe
pas par le point singulier $x$ et qui est lisse sur $S$, et
$U = X - Z$.\footnote{La page dit seulement « nous remplaçons $X$ par
$U = X - Z$ », $Z$ étant « à l'infini ». Les deux conditions sur
$Z$ sont celles qui rendent le calcul légitime, et nous les ajoutons : $Z$
lisse sur $S$ assure que la cohomologie de $U_{\bar\eta}$ se calcule par les
cycles proches sur $U_{s}$, et $x \notin Z$ que les cycles évanescents de
$U$ sont ceux de $X$.} La fibre $U_{s}$ est alors le cône \emph{affine} sur
la quadrique de base, et $U_{\bar\eta}$ une quadrique affine lisse de
dimension $n - 1$. On a la suite\footnote{Les deux derniers termes portent $X$ sur la page, non $U$ ; la
suite n'est exacte qu'avec $U$ partout.}
\[
  \cdots \to H^{i}(U_{s}) \to H^{i}(U_{\bar\eta}) \to H^{i}(X_{s}, R\Phi\Lambda)
  \to H^{i+1}(U_{s}) \to H^{i+1}(U_{\bar\eta}) \to \cdots
\]
et $H^{i}(U_{s}) = 0$ pour
$i \neq 0$, le cône affine étant acyclique (page 10). D'où
$H^{i}(U_{\bar\eta}) \simeq H^{i}(X_{s}, R\Phi\Lambda)$ pour $i \geq 1$, et
\[
  H^{i}(X_{s}, R\Phi\Lambda) = 0 \ \text{pour}\ i \neq n - 1, \qquad
  H^{n-1}(X_{s}, R\Phi\Lambda) \simeq H^{n-1}(U_{\bar\eta}),
\]
groupe libre de rang un engendré par le « vanishing cycle ». Sur les
nombres complexes, la quadrique affine
$x_{0}^{2} + \cdots + x_{n-1}^{2} = 1$ se rétracte sur la sphère réelle de
dimension $n - 1$, et c'est elle qui s'évanouit dans le sommet du cône :
l'image du résumé est ici un théorème.\footnote{L'identification de la
quadrique affine au fibré tangent de la sphère, et de sa fibre à la fibre
de Milnor de la singularité, est notre commentaire.} Le twist,
$\Lambda(-\nu)$ dans les deux parités, est celui que la page 19 calcule
pour $H^{n-1}(U^{n-1})$.

\pagerange{8}{10}
\section*{La cohomologie d'un cône projectif (pages 8 à 16)}

\subsection*{Le cône et le cône épointé (page 10)}

Soit $Y \subset \mathbf{P}^{N}$ une variété projective sur un corps
séparablement clos, $C$ le cône affine sur $Y$, de sommet $a$, et
$\widehat{C}$ le cône projectif, qui contient $C$ comme ouvert et $Y$ comme
section à l'infini, $j : Y \to \widehat{C}$.\footnote{La chemise de la page 8
porte « Calculs de cohomologie des quadriques, des cônes projetants » ;
la page 10 s'ouvre sur « Cohomologie [un mot illisible] d'un cône
projectif ». La page note $X$ la base du cône ; nous la notons $Y$ pour
ne pas la confondre avec l'espace total des pages 2 à 6.} On note
$E = C - a$ et $\widehat{E} = \widehat{C} - a$. L'éclaté de $\widehat{C}$
au sommet est un fibré en droites projectives $\widehat{C}'$ sur $Y$, qui
contient l'éclaté $C'$ de $C$, espace total du fibré
$\mathcal{O}_{Y}(-1)$\footnote{Le diagramme de la page porte aussi $X$ au-dessus de $C'$ et une
flèche $h$ de $X$ vers $C'$, telle qu'elle est tracée, et une flèche $h$
de $X'$ vers $X$. Nous lisons $C'$ comme le fibré en droites au-dessus de
$Y$, donc une projection $C' \to Y$, et $X' \to X$ comme l'identité de la
section à l'infini ; les deux flèches courbes du croquis ne sont pas
reprises.} :

\begin{tikzcd}
C' \arrow[r, hook] \arrow[d] & \widehat{C}' \arrow[d] & Y \arrow[l, hook'] \arrow[d, "\mathrm{id}"] \\
C \arrow[r, hook] & \widehat{C} & Y \arrow[l, hook', "j"']
\end{tikzcd}

La suite de l'éclatement, comparée à celle de $C' \to C$ où
$H^{i}(C') \to H^{i}(Y)$ est bijectif, donne
\[
  H^{i}(C) = 0 \quad (i \neq 0) .
\]
Comme $\widehat{C} - Y = C$ et que $\widehat{C}$ est, au voisinage de $Y$,
un fibré en droites sur $Y$, la suite de Gysin de $j$ donne
\[
  j_{*} : H^{i-2}(Y)(-1) \xrightarrow{\ \sim\ } H^{i}(\widehat{C}) \qquad (i \neq 0),
\]
et la structure multiplicative est déterminée par
$j_{*}(u)\, j_{*}(v) = j_{*}(uv\xi)$, qui vient de la formule de
projection et de $j^{*}j_{*}(w) = w\,\xi$, le fibré normal de $Y$ étant
$\mathcal{O}_{Y}(1)$.\footnote{La page écrit $j_{*}(a)j_{*}(b) = j_{*}(ab\xi)$, la lettre $a$ servant aussi au sommet. Elle fait passer la bijectivité par la
suite exacte courte
$0 \to H^{i}(\widehat{C}) \to H^{i}(\widehat{C}') \to H^{i}(Y) \to 0$ et
écrit « provenant de l'injection de Gysin » (les deux derniers mots douteux). Le passage par
la suite de Gysin de $j$ n'exige pas $Y$ lisse : le couple
$(\widehat{C}, Y)$ est, près de $Y$, celui de la section nulle d'un fibré
en droites.} La classe hyperplane de $\widehat{C}$ est $j_{*}(1)$.

Le cône épointé $\widehat{E}$ est le fibré en droites $\mathcal{O}_{Y}(1)$
au-dessus de $Y$, de section nulle $Y$, donc
$H^{*}(Y) \xrightarrow{\sim} H^{*}(\widehat{E})$, et
$\widehat{E} - Y = E$. La suite de Gysin de $Y \subset \widehat{E}$ s'écrit
\[
  H^{i-2}(Y)(-1) \xrightarrow{\ \xi\ } H^{i}(Y) \to H^{i}(E) \to
  H^{i-1}(Y)(-1) \xrightarrow{\ \xi\ } H^{i+1}(Y),
\]
où la première flèche est le cup-produit par $\xi$, classe d'Euler du
fibré normal. Autrement dit
\[
  0 \to \mathrm{Coker}\bigl(\xi : H^{i-2}(Y)(-1) \to H^{i}(Y)\bigr) \to H^{i}(E)
  \to \mathrm{Ker}\bigl(\xi : H^{i-1}(Y)(-1) \to H^{i+1}(Y)\bigr) \to 0 .
\]

\pagerange{12}{12}
\subsection*{Base lisse, et cohomologie locale au sommet (page 12)}

Supposons $Y$ lisse, connexe, de dimension $d$. Le théorème de Lefschetz
difficile dit que $\xi : H^{j}(Y) \to H^{j+2}(Y)(1)$ est injectif pour
$j \leq d - 1$ et surjectif pour $j \geq d - 1$, à torsion près ; de la
suite précédente on tire, à torsion près,
\[
  \begin{array}{ll}
    i \leq d : & H^{i}(E) \simeq P^{i}(Y) = \mathrm{Coker}\bigl(\xi : H^{i-2}(Y)(-1) \to H^{i}(Y)\bigr), \\[4pt]
    i \geq d + 1 : & H^{i}(E) \simeq \mathrm{Ker}\bigl(\xi : H^{i-1}(Y) \to H^{i+1}(Y)(1)\bigr)(-1)
      \simeq P^{2d+1-i}(Y)(d-i),
  \end{array}
\]
$P$ désignant la cohomologie primitive.\footnote{La page écrit, pour
$i \geq n + 1$, « $H^{i}(E) \simeq P^{i-2}(X)(-1)$ mod.\ de
torsion », l'exposant douteux, la lettre $P$ repassée en gras et son exposant empâté. Le
groupe qui apparaît est un noyau dans $H^{i-1}$, non un conoyau dans
$H^{i-2}$ ; par Lefschetz difficile il a le rang du primitif de degré
$2d + 1 - i$. La marge dit « en vertu du th.\ de Lefschetz ». Il faut
le théorème \emph{difficile} ; en cohomologie $\ell$-adique sur un corps de
caractéristique $p > 0$, il n'est démontré que par Deligne (Weil II, 1980),
et les notes ne pouvaient l'invoquer que sur $\mathbf{C}$, par la théorie
de Hodge, ou pour les quadriques, où on le vérifie à la main, comme aux
pages 14 et 31.} Topologiquement, $E$ se rétracte sur l'entrelacs du sommet,
un fibré en cercles au-dessus de $Y$, et la formule est la suite de Gysin
de ce fibré.

La suite de cohomologie locale au sommet,
\[
  \cdots \to H^{i-1}(C) \to H^{i-1}(E) \to H^{i}_{a}(C) \to H^{i}(C) \to
  H^{i}(E) \to \cdots,
\]
avec $H^{i}(C) = 0$ pour $i \neq 0$ et $H^{0}(C) \to H^{0}(E)$ bijectif
($Y$ connexe), donne
\[
  H^{0}_{a}(C) = H^{1}_{a}(C) = 0, \qquad
  H^{i}_{a}(C) \simeq H^{i-1}(E) \quad (i \geq 2) .
\]

\emph{Le cas d'une quadrique.} Si $Y = Q^{d}$, le cône projectif
$\widehat{C}$ est une quadrique singulière $Q'^{\,d+1}$, et
$H^{i}(Q'^{\,d+1}) \simeq H^{i-2}(Q^{d})(-1)$ pour $i \neq 0$. Avec les
bases des pages 23 et 27, et $\varphi'_{\nu+1} = j_{*}(\varphi_{\nu})$,
$\lambda'_{\nu+1} = j_{*}(\lambda_{\nu})$ :
\[
  \begin{array}{ll}
    d + 1 = 2\nu : & 1, \xi, \ldots, \xi^{\nu}, \ \xi_{\nu+1}, \ldots, \xi_{2\nu} \\[6pt]
    d + 1 = 2\nu + 1 : & 1, \xi, \ldots, \xi^{\nu}, \ [\xi^{\nu+1}, \varphi'_{\nu+1}], \
      \xi_{\nu+2}, \ldots, \xi_{2\nu+1}
  \end{array}
\]
avec $\xi\lambda'_{\nu+1} = 0$ et $\lambda'^{\,2}_{\nu+1} = 0$, par la
formule $j_{*}(u)j_{*}(v) = j_{*}(uv\xi)$ et $\xi\lambda_{\nu} = 0$ sur
$Q^{d}$.\footnote{La page donne la base du second cas sous la forme
$[\xi^{\nu+1}, \lambda'_{\nu+1}]$, qui n'est une base qu'avec $2$
inversible ; voir la page 4. Ce sont, avec un décalage de $\nu$, les bases
de $X_{s}$ écrites à la page 4.}

\pagerange{14}{14}
\subsection*{Le cône épointé sur une quadrique (page 14)}

On applique la suite de la page 10 à $Y = Q^{d}$, en se servant du
cup-produit par $\xi$ de la page 31. Le résultat est la cohomologie de
l'entrelacs d'un point double ordinaire de dimension $d + 1$.

\emph{$d + 1 = 2\nu$, $d = 2\nu - 1$.} Le seul défaut de $\xi$ est en
degré $2\nu - 2 \to 2\nu$, où $\xi^{\nu-1} \mapsto \xi^{\nu} = 2\xi_{\nu}$ ;
le défaut en degré $2d \to 2d + 2$ est la classe du point.\footnote{La page met la torsion en degré $2\nu - 2$, écrit sur
« $2(\nu - 1)$ », avec la liste d'exceptions $0, 2\nu - 2, 4\nu - 1$.
Le conoyau de $\xi$ apparaît dans $H^{i}(E)$ pour $i$ le degré du
\emph{but}, ici $2\nu$. Vérification sur $\nu = 1$ : le cône sur une conique
est le quotient du plan par $\pm 1$, et son entrelacs est l'espace
projectif réel de dimension $3$, de cohomologie entière
$\mathbf{Z}, 0, \mathbf{Z}/2, \mathbf{Z}$. Le groupe $\Lambda/2\Lambda$ est
nul si $m$ est impair, et vaut $\mathbf{Z}/2$ si $m$ est pair, ce que la
page écrit « $0$ si $\ell \neq 2$, $\mathbf{Z}/2\mathbf{Z}$ si
$\ell = 2$ ».}
\[
  H^{i}(E) = 0 \ \ (i \neq 0, 2\nu, 4\nu - 1), \qquad
  H^{0}(E) = \Lambda, \qquad H^{2\nu}(E) \simeq \Lambda/2\Lambda, \qquad
  H^{4\nu-1}(E) \simeq \Lambda(-2\nu) .
\]

\emph{$d + 1 = 2\nu + 1$, $d = 2\nu$.} En degré $2\nu$,
$\xi : \xi^{\nu-1} \mapsto \alpha_{\nu} + \beta_{\nu}$ a pour conoyau
libre l'image de $\varphi_{\nu}$, et
$\xi : H^{2\nu} \to H^{2\nu+2}$ a pour noyau libre $\lambda_{\nu}$ (page 31).\footnote{La page écrit $H^{2\nu-1}(E) \simeq \mu(-\nu+1)$ pour le premier
groupe non trivial, et $2\nu - 1$ dans la liste des exceptions, dont le
$0$ est entouré ; les exposants des trois dernières lignes sont surchargés.
Le conoyau de $\xi$ en degré $2\nu$ donne $H^{2\nu}(E)$, de twist
$-\nu$. Vérification sur $\nu = 1$ : l'entrelacs du cône sur
$\mathbf{P}^{1} \times \mathbf{P}^{1}$ est $S^{2} \times S^{3}$, de
cohomologie en degrés $0, 2, 3, 5$. Les deux autres lignes sont celles de
la page.}
\[
  H^{i}(E) = 0 \ \ (i \neq 0, 2\nu, 2\nu + 1, 4\nu + 1), \qquad
  H^{0}(E) = \Lambda, \qquad H^{2\nu}(E) \simeq \Lambda(-\nu),
\]
\[
  H^{2\nu+1}(E) \simeq \Lambda(-\nu-1), \qquad H^{4\nu+1}(E) \simeq \Lambda(-2\nu-1) .
\]

\pagerange{16}{16}
\subsection*{Supports compacts (page 16)}

Si $Y$ est lisse, $E$ et $\widehat{E}$ sont lisses et leur cohomologie à
supports compacts se déduit de l'ordinaire par dualité de Poincaré ;
$\widehat{C}$ étant propre, ses deux cohomologies coïncident. Reste $C$.
La suite de l'ouvert $C = \widehat{C} - Y$,
\[
  \cdots \to H^{i}(\widehat{C}) \to H^{i}(Y) \to H^{i+1}_{c}(C) \to
  H^{i+1}(\widehat{C}) \to H^{i+1}(Y) \to \cdots,
\]
où $H^{i}(\widehat{C}) \simeq H^{i-2}(Y)(-1)$ et où la restriction
$j^{*}j_{*}$ est le cup-produit par $\xi$, donne pour $i \geq 1$
\[
  0 \to \mathrm{Coker}\bigl(\xi \text{ sur } H^{i-2}(Y)(-1)\bigr) \to
  H^{i+1}_{c}(C) \to \mathrm{Ker}\bigl(\xi \text{ sur } H^{i-1}(Y)(-1)\bigr) \to 0,
\]
c'est-à-dire la même suite que pour $H^{i}(E)$ à la page 10. Les notes en
concluent qu'il faut « prévoir », pour $i \neq 0$,
$H^{i}(E) \simeq H^{i+1}_{c}(C) \simeq H^{i+1}_{a}(C)$, et posent la
question encadrée
\[
  H^{*}_{a}(C) \simeq H^{*}_{c}(C) \ ?
\]
La réponse est oui. Le cône est muni de l'action de $\mathbf{G}_{m}$ par
homothéties, qui contracte $C$ sur son sommet, et pour un faisceau
constant, ou plus généralement monodromique, la cohomologie à supports
compacts de $C$ est la cohomologie à supports dans le sommet.\footnote{C'est
le lemme de contraction, sous la forme que lui ont donnée T.~A.~Springer
(1984) et T.~Braden (2003) ; les notes ne pouvaient pas le connaître.
Topologiquement, $C$ est le cône ouvert sur l'entrelacs $L$, et les deux
groupes valent la cohomologie réduite de $L$ décalée de un. La comparaison
de la page, qui place les deux groupes dans des suites de même forme,
montre seulement que leurs gradués coïncident. Une ligne de la page,
qui compare les deux suites de $\widehat{E} \subset \widehat{C}$ et de
$E$, est presque entièrement illisible.}

\pagerange{19}{19}
\section*{La quadrique lisse et sa partie affine (pages 19 et 21)}

Soit $Q^{n}$ une quadrique lisse de dimension $n \geq 1$ sur un corps $k$
séparablement clos, $Q^{n-1}$ une section hyperplane lisse,
$U^{n} = Q^{n} - Q^{n-1}$, et $A$ un $\Lambda$-module.\footnote{La page note
$U^{n} = AQ^{n}$, « quadrique affine », et dit $A$ « coefficient annulé
par un $m$ », lecture peu sûre.} Alors
\[
  H^{i}_{c}(U^{n}, A) \simeq
  \left\{
  \begin{array}{ll}
    0 & i \neq n, 2n \\
    A \otimes \nu'_{n} & i = n \\
    A(-n) & i = 2n
  \end{array}
  \right.
  \qquad
  H^{i}(U^{n}, A) \simeq
  \left\{
  \begin{array}{ll}
    0 & i \neq 0, n \\
    A & i = 0 \\
    A \otimes \nu_{n} & i = n
  \end{array}
  \right.
\]
avec
\[
  \nu_{n} = \Lambda(-\tfrac{n}{2}), \ \ \nu'_{n} = \Lambda(-\tfrac{n}{2}) \quad (n \text{ pair}),
  \qquad
  \nu_{n} = \Lambda(-\tfrac{n+1}{2}), \ \ \nu'_{n} = \Lambda(-\tfrac{n-1}{2}) \quad (n \text{ impair}),
\]
les deux étant échangés par la dualité de Poincaré :
$\nu'_{n} = \check\nu_{n} \otimes \Lambda(-n)$.\footnote{La page donne
d'abord $\nu'_{n} = \mu_{m}(-\tfrac{n+1}{2})$ pour $n$ impair, puis, d'une
autre encre, $\nu'_{n} = \nu_{n} \otimes \mu_{m}(1)$ pour $n$ impair et
$\nu'_{n} = \check\nu_{n} = \mu_{m}(-\tfrac{n}{2})$ pour $n$ pair. L'ajout
pour $n$ impair est juste et corrige la première ligne ; pour $n$ pair il
faut lire $\nu'_{n} = \nu_{n}$, le dual $\check\nu_{n}$ étant
$\Lambda(\tfrac{n}{2})$ (les accents sont peu lisibles). Pour $\nu_{n}$, $n$
impair, le signe de $n \pm 1$ est surchargé ; c'est $n + 1$. Dans le
tableau de $H^{i}(U^{n}, A)$, deux flèches croisées sur la page échangent
les deux dernières lignes, et la lecture rectifiée est celle donnée ici. Le
cas $n = 0$, où $U^{0} = Q^{0}$ a deux points, est mis à part par un
« si $n \neq 0$ » entouré.} Ces formules se déduisent de la suite de
Gysin
\[
  \cdots \to H^{i-1}(U^{n}) \to H^{i-2}(Q^{n-1})(-1) \xrightarrow{\ \alpha_{*}\ }
  H^{i}(Q^{n}) \to H^{i}(U^{n}) \to H^{i-1}(Q^{n-1})(-1) \to \cdots
\]
et de la cohomologie des quadriques lisses (pages 23 à 27) : $\alpha_{*}$
est surjectif pour $i \neq 0, n$, injectif pour $i \neq 1, n + 1$, et l'on
calcule séparément les degrés restants. Pour $i = 0, 1$ la source est nulle.
Pour $i = n, n + 1$, il reste
\[
  0 \to H^{n-2}(Q^{n-1})(-1) \to H^{n}(Q^{n}) \to H^{n}(U^{n}) \to
  H^{n-1}(Q^{n-1})(-1) \to H^{n+1}(Q^{n}) \to 0 .
\]
Si $n$ est pair, les deux derniers termes sont nuls, et $H^{n}(U^{n})$ est
le quotient de $H^{n}(Q^{n}) = \Lambda\alpha \oplus \Lambda\beta$ par
l'image $\Lambda(\alpha + \beta)$ de $\xi^{n/2-1}$ : il est engendré par
l'image $\ell^{n}$ de $\alpha$, celle de $\beta$ valant $-\ell^{n}$. Si $n$ est
impair (page 21), ce sont les deux premiers, et $H^{n}(U^{n})$ est le noyau
de $H^{n-1}(Q^{n-1})(-1) \to H^{n+1}(Q^{n})$, qui envoie les classes des deux
familles de $Q^{n-1}$ sur la classe d'un même sous-espace linéaire : il est
engendré par $\lambda_{(n-1)/2}$. Dans les deux cas, le générateur dépend du
choix d'une famille, et l'isomorphisme n'est déterminé qu'au signe
près.\footnote{La marge de la page 19, écrite en biais à côté d'une double
flèche qui relie $H^{i}_{c}$ et $H^{i}$, dit « si $n$ impair, [deux mots
illisibles], on ne sait déterminer qu'au signe près ». L'ambiguïté vaut
aussi pour $n$ pair, comme on vient de le voir.}

\emph{L'application d'oubli des supports}
$H^{i}_{c}(U^{n}, A) \to H^{i}(U^{n}, A)$ se factorise par
$H^{i}(Q^{n}, A)$ ; elle est donc nulle, sauf pour $n$ pair et $i = n$.
Dans ce cas $H^{n}_{c}(U^{n})$ est engendré par un élément $\ell_{n}$
d'image $\lambda_{n/2}$ dans $H^{n}(Q^{n})$, et son image dans
$H^{n}(U^{n})$ est $2\ell^{n}$ ; l'accouplement de Poincaré donne\footnote{La page écrit $\langle \ell, \varphi(\ell) \rangle = 2$, précédé
d'une ligne entière biffée et d'un passage illisible où se lit
« dans exc.~2 », lecture douteuse. Le signe est celui de l'auto-intersection
du cycle évanescent, $(-1)^{n/2}\,2$ en dimension paire : voir le carré de
$\lambda_{\nu}$ à la page 25.}
\[
  \langle \ell_{n}, 2\ell^{n} \rangle = \lambda_{n/2}^{2} = (-1)^{n/2}\, 2 .
\]

\pagerange{21}{21}
\subsection*{Le programme de la page 21}

La page résume ce qui précède et ce qui reste à établir, pour $Q^{n}$
lisse\footnote{Les énoncés c) et c') de la page sont écrits vite et ne sont pas
justes tels quels : c) dit la restriction « bijectif si $i \geq 2$,
$i \neq n$, injectif pour tt $i$ » ($n$ pair) et « bijectif si
$i \geq 2$, $i \neq n + 1$, surjectif pour tt $i \geq 1$ » ($n$ impair) ;
c') dit l'application dans l'autre sens « bijectif si $i \leq 2n - 2$,
$i \neq n$, injectif pour tt $i$ » ($n$ pair) et « bijectif si
$i \leq 2n - 2$, $i \neq n - 1$, injectif pour tt $i \leq 2n - \ldots$ »
($n$ impair), avec plusieurs exposants douteux et un mot entouré illisible.
Nous donnons les énoncés exacts, qui se lisent sur les bases des pages 23
et 27 ; la restriction n'est pas injective en degré $2n + 2$, où le but est
nul, et l'application de Gysin n'est pas injective en degré moyen pour $n$
pair. Pour d), la page écrit « $i = n$ » : c'est le degré $2i = n$.} :

\begin{enumerate}
\item[a)] la cohomologie de $Q^{n}$ est nulle en degré impair ;
\item[b)] $H^{i}(U^{n})$ est nul pour $i \neq 0, n$, isomorphe à $A$ pour
$i = 0$ et à $A \otimes \nu_{n}$ pour $i = n$ ; b') l'énoncé dual, pour les
supports compacts ;
\item[c)] la restriction $H^{i}(Q^{n+1}) \to H^{i}(Q^{n})$ est bijective
pour $i \leq 2n$, sauf en un degré : si $n$ est pair, elle est injective de
conoyau libre de rang un pour $i = n$ ; si $n$ est impair, elle est
surjective de noyau libre de rang un pour $i = n + 1$ ;
\item[c')] dualement, l'application de Gysin
$H^{i}(Q^{n})(-1) \to H^{i+2}(Q^{n+1})$ est bijective pour $i \geq 0$, sauf
en un degré : si $n$ est pair, elle est surjective de noyau
$\Lambda\lambda_{n/2}$ pour $i = n$ ; si $n$ est impair, elle est injective
de conoyau libre de rang un pour $i = n - 1$ ;
\item[d)] $H^{2i}(Q^{n})$ est libre de rang $1$ pour $0 \leq i \leq n$, sauf
pour $n$ pair et $2i = n$, où il est de rang $2$.
\end{enumerate}

Le
point e) commence par « si $n$ pair » et le reste est
illisible.\footnote{Une ligne au crayon, tête-bêche au pied de la page,
« $X_{1} \subset X' = Y[t_{1}, \ldots, t_{n}]$, les deux s'envoyant sur
$Y$ », ne se rattache à rien de ce qui précède.}

\pagerange{23}{27}
\section*{L'anneau de cohomologie des quadriques (pages 23 à 34)}

\subsection*{Bases (pages 23 et 27)}

Pour une quadrique lisse $Q^{n}$, $\xi^{i} \in H^{2i}(Q^{n}, \Lambda(i))$, et
pour $n = 2\nu$ la classe supplémentaire $\lambda_{\nu}$ en degré moyen. Les
bases entières de la cohomologie sont\footnote{Le cadre de la page écrit le degré moyen de $Q^{2\nu}$ sous la
forme $[\xi^{\nu}, \lambda_{\nu}]$ ; une note en biais au-dessus du cadre,
reliée par un trait à ce crochet, introduit
$\tfrac12(\xi^{\nu} + \lambda_{\nu}) = \varphi_{\nu}$ et la base
$(\xi^{\nu}, \varphi_{\nu})$, à cause d'un « $2$ dans
$H^{2\nu}(Q^{2\nu})$ ». C'est la base entière, et nous la prenons partout.
La matrice $\left(\begin{smallmatrix} 1 & 1 \\ 1 & -1 \end{smallmatrix}\right)$
des pages 23 et 27 est celle du passage de $(\alpha_{\nu}, \beta_{\nu})$ à
$(\xi^{\nu}, \lambda_{\nu})$ ; son déterminant vaut $-2$.}
\[
  \begin{array}{ll}
    Q^{2\nu} : & \xi^{0}, \ldots, \xi^{\nu-1}, \ [\xi^{\nu}, \varphi_{\nu}], \
      \xi_{\nu+1}, \ldots, \xi_{2\nu} \\[4pt]
    Q^{2\nu+1} : & \xi^{0}, \ldots, \xi^{\nu}, \ \xi_{\nu+1}, \ldots, \xi_{2\nu+1}
  \end{array}
\]
 La page 23 dresse
le tableau pour $Q^{0} \subset Q^{1} \subset \cdots \subset Q^{4}$, la page
27, écrite en travers, le prolonge jusqu'à $Q^{6}$\footnote{Les crochets $[\lambda_{\nu}, \xi^{\nu}]$ sont ceux de la page, à
lire avec la réserve de la note précédente. Pour $Q^{0}$, deux points,
$\xi^{0}$ est la somme des deux points et $\lambda_{0}$ leur différence ;
la page écrit cette colonne « $\lambda_{0}^{+}\,\xi^{0}_{+}$ », indices
douteux. Elle commence par «
$\mathbf{P}^{1} \times \mathbf{P}^{1} \simeq Q^{2} \subset Q^{1} \simeq
\mathbf{P}^{1}$ », l'inclusion dans le mauvais sens : c'est $Q^{1}$ qui
est une section de $Q^{2}$. Au pied de la page 27, tête-bêche et biffé, se
lit le a) de la page 21.} :
\[
  \begin{array}{ccccccc}
    Q^{0} & Q^{1} & Q^{2} & Q^{3} & Q^{4} & Q^{5} & Q^{6} \\[4pt]
    [\lambda_{0}, \xi^{0}] & 1 & 1 & 1 & 1 & 1 & 1 \\
     & \xi_{1} & [\lambda_{1}, \xi] & \xi & \xi & \xi & \xi \\
     & & \xi_{2} & \xi_{2} & [\lambda_{2}, \xi^{2}] & \xi^{2} & \xi^{2} \\
     & & & \xi_{3} & \xi_{3} & \xi_{3} & [\lambda_{3}, \xi^{3}] \\
     & & & & \xi_{4} & \xi_{4} & \xi_{4} \\
     & & & & & \xi_{5} & \xi_{5} \\
     & & & & & & \xi_{6}
  \end{array}
\]

\pagerange{25}{25}
\subsection*{Relations, restriction et Gysin (pages 23 et 25)}

Dans $Q^{2\nu}$ :
\[
  2\xi_{i} = \xi^{i}, \qquad \xi^{i}\xi_{j} = \xi_{i+j}, \qquad
  \xi\lambda_{\nu} = 0, \qquad
  \lambda_{\nu}^{2} = (-1)^{\nu}\, 2\xi_{2\nu}, \qquad
  (\xi^{\nu})^{2} = \xi^{2\nu} = 2\xi_{2\nu},
\]
les produits $\xi_{i}\xi_{j}$ et $\lambda_{\nu}\xi_{j}$ ($i, j \geq \nu + 1$)
étant nuls pour des raisons de degré ; $\xi_{2\nu}$ est la classe du point.
En effet $\xi\alpha_{\nu} = \xi\beta_{\nu} = \xi_{\nu+1}$, et deux
$\nu$-plans génériques de $Q^{2\nu}$ se coupent en un point s'ils sont de
même famille et $\nu$ est pair, ou de familles opposées et $\nu$ est
impair, et ne se coupent pas sinon.\footnote{La page écrit
$\lambda_{\nu}^{2} = 2\xi_{2\nu} = \xi^{2\nu}$, ce qui n'est vrai que pour
$\nu$ pair. Pour $\nu = 1$, $Q^{2} = \mathbf{P}^{1} \times \mathbf{P}^{1}$
et $\lambda_{1}$ est la différence des deux règles, de carré $-2$. Une
présentation de l'anneau, $\mathbf{Z}[T][T']/(T^{2\nu+2}, 2T' - T^{\nu+1})$,
est biffée en tête de page ; elle vise le cas impair.}

Pour $Q^{i} \subset Q^{j}$, $i \leq j$, section linéaire, la restriction et
l'application de Gysin se lisent sur ces bases :
\[
  \begin{array}{ll}
    j = 2\nu > i : &
      \xi^{\alpha} \mapsto \xi^{\alpha}, \ \ \xi_{\alpha} \mapsto \xi_{\alpha}, \ \
      \lambda_{\nu} \mapsto 0, \ \ \varphi_{\nu} \mapsto \xi_{\nu}
      \quad (\text{restriction}) \\[4pt]
    i = 2\nu < j : &
      \xi^{\alpha} \mapsto \xi^{\alpha+(j-i)}, \ \ \xi_{\alpha} \mapsto \xi_{\alpha+(j-i)}, \ \
      \lambda_{\nu} \mapsto 0, \ \ \varphi_{\nu} \mapsto \xi_{\nu+(j-i)}
      \quad (\text{Gysin})
  \end{array}
\]
La classe $\lambda_{\nu}$ est tuée dans les deux sens. Par l'application de
Gysin, c'est une classe « évanescente » : par la théorie de Lefschetz,
le noyau de Gysin d'une section hyperplane est engendré par les cycles
évanescents du pinceau, et $\lambda_{\nu}$ est la classe, dans la quadrique lisse, du cycle
évanescent des pages 4 et 6.\footnote{La page qualifie $\lambda_{\nu}$ de classe
« effective » dans le cas de la restriction et de classe
« évanescente » dans celui de Gysin, les deux mots douteux
et le verbe illisible. Le commentaire sur les pinceaux de Lefschetz est le
nôtre ; l'adjectif « effective », qui conviendrait plutôt à
$\varphi_{\nu}$, classe d'un plan, n'a pas de lecture sûre. La page 23
donne la règle de Gysin pour $Q^{n-1} \subset Q^{n}$ sous la forme
$\xi_{i} \rightsquigarrow \xi_{i+2}$ ; l'indice monte de $1$, comme le dit
la page 25.}

\pagerange{29}{29}
\subsection*{Le cycle évanescent et la classe $\lambda_{\nu}$ (page 29)}

Soit $n = 2\nu$. Deux suites relient la cohomologie moyenne de $Q^{n}$ aux
quadriques affines voisines, $U^{n} = Q^{n} - Q^{n-1}$ et
$U^{n+1} = Q^{n+1} - Q^{n}$. On note $\ell^{i}$ un générateur de
$H^{i}(U^{i})$ et $\ell_{i}$ un générateur de $H^{i}_{c}(U^{i})$ (page 19).
\[
  H^{n+1}(U^{n+1})(1) \xrightarrow{\ \partial\ } H^{n}(Q^{n})
    \xrightarrow{\ \text{res}\ } H^{n}(U^{n}),
  \qquad
  \ell^{2\nu+1} \mapsto \lambda_{\nu} \mapsto 2\ell^{2\nu}, \quad
  \xi^{\nu} \mapsto 0, \quad \varphi_{\nu} \mapsto \ell^{2\nu} ;
\]
\[
  H^{n}_{c}(U^{n}) \longrightarrow H^{n}(Q^{n}) \xrightarrow{\ \partial\ }
    H^{n+1}_{c}(U^{n+1}),
  \qquad
  \ell_{2\nu} \mapsto \lambda_{\nu} \mapsto 2\ell_{2\nu+1}, \quad
  \xi^{\nu} \mapsto 0, \quad \varphi_{\nu} \mapsto \ell_{2\nu+1} .
\]
La première flèche du haut est le résidu de la suite de Gysin de
$Q^{n} \subset Q^{n+1}$, d'image le noyau $\Lambda\lambda_{\nu}$ de Gysin ;
la seconde du bas est le cobord de la suite de $U^{n+1} \subset Q^{n+1}$,
qui identifie $H^{n+1}_{c}(U^{n+1})$ au conoyau de la restriction
$H^{n}(Q^{n+1}) \to H^{n}(Q^{n})$. Dans les deux suites, le composé est la
multiplication par $2$ entre générateurs : c'est le $2$ de
l'auto-intersection du cycle évanescent.\footnote{La page écrit le premier
terme $H^{n-1}(U^{n-1})(1)$ ; avec ce terme, aucun twist ne s'accorde avec
$H^{n}(Q^{n})$, et la flèche ne peut atteindre $\lambda_{\nu}$. Nous lisons
$H^{n+1}(U^{n+1})(1)$, que confirment la notation $\ell^{2\nu+1}$ de la page
et la ligne $Q^{n+1} - Q^{n} = U^{n+1}$ qu'elle écrit au-dessous. Elle
écrit aussi $H^{n+1}_{c}(U^{n+1})(+1)$ ; sans twist, les deux membres ont
déjà le même, $\Lambda(-\nu)$. L'étiquette « ins » de la flèche du
haut est douteuse, celle de la flèche du bas illisible. Une flèche courbe
du premier terme au troisième porte « bij » et la fraction
« $\mathrm{exc}^{2}/\ell^{(2\nu+1)} \rightsquigarrow 2\ell^{(2\nu)}$ »,
lecture douteuse : le composé n'est bijectif que si $2$ est inversible.
Deux lignes au crayon en tête de page,
$L_{\cdot}(F_{\cdot}) \otimes L_{\cdot}(G_{\cdot}) = F \otimes^{L} G$, sur le
produit tensoriel dérivé, ne se rattachent pas au reste.}

\emph{N.B.} La page note enfin qu'on peut définir $\varphi_{\nu}$
directement par un sous-espace isotrope maximal : les deux familles de
tels sous-espaces donnent les classes $\alpha_{\nu}$, $\beta_{\nu}$, avec
$\alpha_{\nu} + \beta_{\nu} = \xi^{\nu}$,
$\alpha_{\nu} - \beta_{\nu} = \lambda_{\nu}$ et
$\varphi_{\nu} = \alpha_{\nu}$. Ce sont les définitions que cette lecture
a prises dès le départ.\footnote{« directement », « isotrope » et
« maximal » sont des lectures douteuses, et quelques mots sont
illisibles.}

\pagerange{31}{31}
\subsection*{Le cup-produit par $\xi$ (page 31)}

Sur $Q^{2\nu}$, $\xi : H^{2i}(Q^{2\nu}) \to H^{2i+2}(Q^{2\nu})(1)$,
$0 \leq i \leq 2\nu - 1$, est bijectif pour $i \neq \nu - 1, \nu$ ; pour
$i = \nu - 1$ il est injectif, de conoyau libre de base $\varphi_{\nu}$ ;
pour $i = \nu$ il est surjectif, de noyau libre de base $\lambda_{\nu}$.

Sur $Q^{2\nu+1}$, $\xi : H^{2i}(Q^{2\nu+1}) \to H^{2i+2}(Q^{2\nu+1})(1)$,
$0 \leq i \leq 2\nu$, est bijectif pour $i \neq \nu$ ; pour $i = \nu$ il est
injectif de conoyau $\Lambda/2\Lambda$, puisque $\xi^{\nu} \mapsto
\xi^{\nu+1} = 2\xi_{\nu+1}$.\footnote{La page écrit le conoyau
$\mathbf{Z}/2\mathbf{Z}$ ; avec des coefficients $\Lambda$, c'est
$\Lambda/2\Lambda$. Une première ligne, « injectif si $i < 2\nu$ i.e.\
$i \leq \nu - 1$ », est biffée.} C'est la forme du théorème de Lefschetz
difficile pour les quadriques, et c'est ce qui sert à la page 14 :
l'énoncé tient sur les bases entières, au défaut de torsion près en
dimension impaire.

\pagerange{34}{34}
\subsection*{Les groupes de Chow (page 34)}

Une feuille de papier bistre, d'une écriture plus posée, donne les mêmes
bases pour l'équivalence linéaire des cycles. Sur $Q^{2n}$ :
\[
  1, H, H^{2}, \ldots, H^{n-1}, \ S', S'', \ HS' = HS'', \ \ldots, \
  H^{n}S' = H^{n}S'', \qquad H^{n} = S' + S'',
\]
où $H$ est la section hyperplane et $S'$, $S''$ les classes de deux
sous-espaces linéaires de dimension $n$ contenus dans $Q^{2n}$, pris dans
les deux familles algébriques distinctes de tels sous-espaces. Sur
$Q^{2n+1}$ :
\[
  1, H, \ldots, H^{n}, \ \tfrac{H^{n+1}}{2}, \ldots, \tfrac{H^{2n+1}}{2},
\]
les classes $\tfrac{H^{k}}{2}$ ($k \geq n+1$) étant celles des sous-variétés
linéaires contenues dans $Q^{2n+1}$.\footnote{Ici l'indice $n$ est la
demi-dimension, comme sur la page. Le mot sous l'accolade de $S'$, $S''$
est illisible, et « de dim. » une lecture douteuse.} C'est l'anneau
de Chow des quadriques lisses, et l'application classe de cycles l'envoie
isomorphiquement sur la cohomologie paire des pages précédentes.

\section*{Les feuillets du tapuscrit de l'exposé XII de SGA~3}

Quatorze feuillets portent au recto un état dactylographié, corrigé à
l'encre de sa main, de l'exposé « XII Tores maximaux, groupe de Weyl,
sous-groupes de Cartan, centre réductif des schémas en groupes lisses et
affines », sans rapport avec les notes écrites au verso. Le texte tapé
est le sien et publié ; on ne le recompose pas, on en donne la teneur
et ses interventions principales, détaillées dans la
transcription.\footnote{La matière correspond à l'exposé XII de SGA~3
publié, à notre connaissance ; on ne compare pas numéro à numéro. Ses
corrections renumérotent partout les paragraphes en 1.x, 2.x, 4.x, et les
feuillets ne sont pas rangés dans l'ordre du tapuscrit : les pages tapées
lisibles sont 0, 1, 11, 12, 21, 23, 26 et 31. Les notes du verso sont donc
postérieures à cette frappe, qui appartient au séminaire de 1962-1964.}

\pagerange{7}{7}
\subsection*{Page de titre (page 7, tapuscrit p.~0)}

L'exposé annonce qu'il utilisera la théorie de Borel des groupes
algébriques lisses sur un corps algébriquement clos, exposée au Séminaire
Chevalley 1956-1957 (« BIBLE »), la théorie des schémas n'y
apportant pas de simplification notable, pour en tirer des résultats
analogues sur une base quelconque. Il y ajoute de sa main un paragraphe :
les exposés oraux se limitaient aux préschémas en groupes affines, en
s'appuyant sur les théorèmes de représentabilité de l'exposé XI, n° 4 ; les
présentes notes (n° 6 à 8) montrent comment éliminer l'hypothèse affine
par une méthode plus simple, qui n'utilise pas les résultats les plus
délicats de l'exposé XI. En marge : « Il est évident que le contenu des
exposés XI, XII, XV, XVI devrait être complètement refondu ».

\pagerange{5}{5}
\subsection*{Tores maximaux (page 5, tapuscrit p.~1)}

Début du n° 1, « Tores maximaux » (il biffe
« et groupe de Weyl ») : tores maximaux d'un groupe algébrique $G$ sur
$k$ algébriquement clos, qui sont ceux de $G_{\text{réd}}$ ; conjugaison
par $G(k) = G_{\text{réd}}(k)$ (il ajoute l'égalité) ; rang réductif ;
l'hypothèse affine est inutile grâce au théorème de Chevalley sur les
extensions de variétés abéliennes par des groupes affines (il renvoie au
n° 6, et biffe « nous n'aurons pas besoin de cette généralisation ») ;
début de la définition du sous-groupe de Cartan, où il remplace « $G$
sur $k$ étant comme ci-dessus » par « Soit $G$ un groupe algébrique
lisse sur $k$ ».

\pagerange{17}{17}
\subsection*{Les rangs, et l'énoncé du théorème 1.7 (page 17)}

Définition 1.5 des rangs réductif, nilpotent et
unipotent sur un corps, invariants par extension du corps ; Remarque 1.6,
un schéma en groupes affine et lisse sur un trait, de fibre générique
$\mathbf{G}_{m}$ et de fibre spéciale $\mathbf{G}_{a}$, sans tore maximal,
où le rang réductif des fibres n'est pas continu ; énoncé du
Théorème 1.7 sur les fonctions $\rho_{r}$, $\rho_{n}$,
$\rho_{u} = \rho_{n} - \rho_{r}$. Dans la remarque qui précède, sur les
tores contenus dans un tore maximal, il ajoute qu'on verra dans l'exposé XIV
que la réciproque est vraie lorsque $S$ est artinien, ou local avec $G$
réductif.

\pagerange{3}{3}
\subsection*{Le théorème 1.7 (page 3)}

Suite de 1.7 : a) $\rho_{r}$ est semi-continue
inférieurement, $\rho_{n}$ et $\rho_{u}$ supérieurement ; b) $\rho_{r}$
est localement constante si et seulement si $G$ admet localement, pour la
topologie étale (ou fidèlement plate quasi-compacte), un tore maximal ;
c) deux tores maximaux sont localement conjugués pour la topologie
étale ; d) sous b), $\rho_{n}$ et $\rho_{u}$ sont localement constantes.
Il insère « de Zariski » dans la réduction au cas d'un changement de
base fidèlement plat.

\pagerange{13}{13}
\subsection*{Fin de 1.7, corollaires 1.9 et 1.10 (page 13)}

Fin de la preuve de 1.7 ; Corollaire 1.9 (si
$\rho_{u}(s) = 0$, les rangs sont constants au voisinage de $s$) ;
Corollaire 1.10 (pour $G$ de rang réductif localement constant, le
foncteur $\mathcal{T}$ des tores maximaux est représentable par un
préschéma lisse, séparé, de type fini, isomorphe à $G/\mathrm{Norm}_{G}(T)$
quand il existe un tore maximal $T$). Il ajoute les renvois à l'exposé XI.

\pagerange{11}{11}
\subsection*{Le préschéma des tores maximaux (page 11)}

Remarques 1.11 : $\mathcal{T}$ est le \emph{préschéma des
tores maximaux}, affine sur $S$ (n° 5). Il remplace la fin tapée, encadrée
et barrée, par : « de façon générale, il est possible (et sans doute
avantageux) d'éliminer totalement les résultats de représentabilité de
Exp XI N° 4 de la théorie présentée dans le présent exposé et les
suivants ».

\pagerange{24}{24}
\subsection*{Tores centraux (page 24, tapuscrit p.~11)}

Lemme 1.14 : un sous-tore qui commute à un tore maximal y
est contenu ; Corollaires 1.15 et 1.16 : un groupe commutatif, ou à fibres
connexes nilpotentes, lisse et affine, de rang réductif localement
constant, a un unique tore maximal, central. Il remplace « affine »
par « un $S$-préschéma en groupes » dans le lemme et insère
l'hypothèse de rang réductif localement constant dans les deux corollaires.

\pagerange{28}{28}
\subsection*{Le Nullstellensatz, et le début du groupe de Weyl (page 28, tapuscrit p.~12)}

Fin de 1.16 : il remplace un renvoi à l'exposé VI par un
argument écrit en marge, que $\mathrm{Centr}_{G}(T)$ est un sous-schéma
fermé contenant $G(k)$, donc égal à $G$ par le Nullstellensatz, $G$ étant
réduit. Proposition 1.17, qu'il réécrit en équivalence : pour
$G \supset H \supset T$, $T$ est un tore maximal de $G$ si et seulement si
c'est un tore maximal de $H$ et si les rangs réductifs des fibres de $G$
et de $H$ coïncident. Début du n° 2, « Le groupe de Weyl » :
$W = N/C$, groupe fini étale, indépendant du tore maximal par le théorème
de conjugaison de Borel.

\pagerange{20}{20}
\subsection*{Le groupe de Weyl : lemme 2.2 (page 20)}

Lemme 2.2 : $W$ étale, séparé, quasi-fini sur $S$ ; la
fonction qui associe à $s$ la classe de la fibre géométrique de $W$ est
semi-continue inférieurement, et constante au voisinage de $s$ si et
seulement si $W$ est fini au-dessus d'un voisinage ; esquisse de la
démonstration.

\pagerange{15}{15}
\subsection*{Lemme 2.3 (page 15)}

Fin de cette démonstration, par le critère valuatif de
propreté ; Lemme 2.3 : pour $R \subset T$ deux sous-tores d'un groupe
affine lisse sur $k$ algébriquement clos, $W(R)$ est quotient d'un
sous-groupe de $W(T)$. Il précise « fini sur $S$ » et
« monomorphisme », et barre le dernier alinéa.

\pagerange{26}{26}
\subsection*{Centre réductif (page 26, tapuscrit p.~21)}

Critère 4.3 : $Z$ est un centre réductif de $G$ si et
seulement s'il est de type multiplicatif et si chaque $Z_{s}$ est un centre
réductif de $G_{s}$ ; Théorème 4.4 : un groupe algébrique affine sur un
corps admet un centre réductif. Il ajoute « à fibres affines » et le
renvoi à l'exposé XIII pour la décomposition $G = Z \times U$.

\pagerange{22}{22}
\subsection*{Le quotient par le centre réductif (page 22, tapuscrit p.~23)}

Fin de la preuve que $G/Z$ n'a pas de sous-groupe de type
multiplicatif non trivial, et début de la preuve que $Z$ est un centre
réductif ; il remplace un mot par « implique ».

\pagerange{30}{30}
\subsection*{Fin de la démonstration (page 30, tapuscrit p.~26)}

Démonstration de ce que $G/Z$ a pour centre réductif le
groupe unité ; de sa main, un seul crochet. Le feuillet est gardé pour ne
pas rompre la suite.

\pagerange{32}{32}
\subsection*{Un dernier feuillet (page 32, tapuscrit p.~31)}

Démonstration en quatre points d'un énoncé sur le rang réductif, parallèle
au Théorème 1.7 : un tore maximal dans une fibre le reste au voisinage si
$\rho_{r}$ est localement constant ; le foncteur des tores maximaux,
sous-foncteur ouvert du foncteur $F$ de l'exposé XI, 4.1, est lisse et
séparé et admet des sections localement pour la topologie étale. Il
complète les hypothèses (« et séparé », « moyennant (i) ») et les
renvois.

\pagerange{35}{37}
\section*{Le tiré à part (pages 35 à 38)}

Le dossier se termine sur le tiré à part d'une Note aux Comptes rendus :
Antoine Delzant, « Définition des classes de Stiefel-Whitney d'un module
quadratique sur un corps de caractéristique différente de 2 »,
\emph{C. R. Acad. Sci. Paris} 255 (1962), p.~1366-1368 ; il ne porte rien
de sa main. La Note construit, à l'aide de l'anneau filtré
$\widetilde{H}$ que Grothendieck associe à un anneau gradué $H$ (Bull.\ SMF
86, 1958), un homomorphisme de l'anneau des classes de formes
quadratiques dans $\widetilde{H}$ pour $H$ la cohomologie galoisienne à
coefficients $\mathbf{Z}/2$, d'où des classes
$w^{i}(q) \in H^{i}(k, \mathbf{Z}/2)$ vérifiant la formule de Whitney ;
$w^{1}$ est le discriminant, $w^{2}$ l'invariant de Hasse-Witt. Sur un corps
de nombres, une forme est déterminée par son rang et ses classes. La Note
se ferme sur deux questions : l'application du gradué de degré $2$ vers
$H^{2}(k, \mathbf{Z}/2)$ est-elle injective, et tout élément d'ordre $2$
du groupe de Brauer est-il somme de classes d'algèbres de quaternions
?\footnote{Les deux réponses sont oui, par le théorème de Merkurjev
(1981) ; leur généralisation en tout degré est la conjecture de Milnor,
démontrée par Voevodsky et Orlov-Vishik-Voevodsky. La page 38, dos du
tiré à part, ne porte que la mention imprimée de son origine. Rien ne
relie le tiré à part aux notes, sinon les quadriques.}

\end{document}
